\cotefile{60}
\scanpage{48}
\mspage{4 - 41 -}
\noindent\hfill n\textsuperscript{o} 392\par

principal semi-local, dont les idéaux maximaux sont les $\underline{p}S^{-1}A$, où $\underline{p}\in Y$, les anneaux locaux correspondants étant les $A_{\underline{p}}$. On a donc $S^{-1}A/\underline{p}S^{-1}A=k(\underline{p})$, d'où $S^{-1}M/\underline{p}S^{-1}M=(M/\underline{p}M)\otimes_A S^{-1}A=M\otimes_A(A/\underline{p}\otimes_A S^{-1}A)=M\otimes_A k(\underline{p})=M(\underline{p})$ pour tout $\underline{p}\in Y$. En vertu du chap. II, \S~1, n\textsuperscript{o} 2, prop. 6, il existe un élément $z/s\in S^{-1}M$, $(z\in M,\ s\in S)$, dont les images canoniques dans les $M(\underline{p})$ pour $\underline{p}\in Y$ sont toutes $\ne 0$. Par définition de $S$, on a donc $z(\underline{p})\ne 0$ pour tout $\underline{p}\in Y$. On peut en outre supposer que $y$ et $z$ sont \uline{linéairement indépendants} sur $K$. En effet, dans le cas contraire, considérons un élément $t\in M$ linéairement indépendant de $y$, (il en existe puisque $n\geq 2$) ; prenons d'autre part un élément $a\ne 0$ appartenant à $\handwritten{\bigcap}_{\underline{p}\in Y}\underline{p}$, (qui n'est pas réduit à 0 puisque $A$ est intègre), et posons $z'=z+at$ : il est clair que $y$ et $z'$ sont linéairement indépendants sur $K$ et que l'on a $z'(\underline{p})=z(\underline{p})\ne 0$ pour tout $\underline{p}\in Y$.

Supposant donc $y$ et $z$ linéairement indépendants sur $K$, soit $Z$ l'ensemble des $\underline{p}\in\unsure{,} P-Y$ tels que $y(\underline{p})$ et $z(\underline{p})$ soient linéairement dépendants sur $k(\underline{p})$ ; il résulte du lemme 7 que cet ensemble est \uline{fini}. Pour tout $\underline{p}\in Z$, on peut donc écrire $z(\underline{p})=\lambda(\underline{p})y(\underline{p})$ avec $\lambda(\underline{p})\in k(\underline{p})$. Or, on a $\operatorname{Card}(A/\underline{p})\geq 2$ pour tout $\underline{p}\in P$ ; il résulte donc du lemme 6 qu'il existe $b\in A$ tel que, pour tout $\underline{p}\in Z$, l'image canonique de $b$ dans $A/\underline{p}$ soit distincte de $\lambda(\underline{p})$. Montrons alors que l'élément $x=z-by$ répond à la question ; il suffit, (en vertu du lemme 7 appliqué pour $m=1$), de vérifier que $x(\underline{p})\ne 0$ pour \uline{tout} $\underline{p}\in P$. Or :

- si $\underline{p}\in Y$, on a $x(\underline{p})=z(\underline{p})\ne 0$ par construction ;

- si $\underline{p}\in Z$,\unsure{\textbackslash} on a $x(\underline{p})=\mu_{\underline{p}}\unsure{\circ|\cdot}y(\underline{p})$ avec $\mu_{\underline{p}}\ne 0$ en vertu du choix de $b$, donc $x(\underline{p})\ne 0$ puisque $y(\underline{p})\ne 0$ ;

- si $\underline{p}\in P-(Y\cup Z)$, $y(\underline{p})$ et $z(\underline{p})$ sont linéairement indépendants, donc $x(\underline{p})\ne 0$.

Ces lemmes étant établis, passons à la démonstration du th. 6. Raisonnons par récurrence sur $n$, le cas $n\leq 1$ étant trivial puisque $M$ lui-même est alors iso-
% archive sheet number: 47
\scanpage{49}
Calcul de $D(E_\eta^{\unsure{\leftrightharpoons|\leftrightarrows|\natural}})$ :

\[
\begin{aligned}
D(E_\eta^{\leftrightharpoons}) &= D(i_!(E_U))+\sum_{x\in Y}D(j_{x*}(E_\eta^{\leftrightharpoons(x)}))\\
&=\mathbb R i_*(D E_U)+\sum_{x\in Y}j_{x*}(D(E_\eta^{\leftrightharpoons(x)}))
\end{aligned}
\]

Or la formation de $\leftrightharpoons$ local \unsure{commute} au \unsure{passage} au dual, d'autre part on a
\[
D(E_U)=\check E_U(1)\deleted{[2]E_U(\rho+1)\unsure{=|,}},
\]
d'où

\[
\begin{aligned}
D(E_\eta^{\leftrightharpoons})&=(\mathbb R i_*(\check E_U))(\deleted{\rho+}1)+\sum_{x\in Y}j_{x*}(\check E_\eta^{\leftrightharpoons(x)})\\
&\deleted{=}\qquad\deleted{\unsure{i_*|i}}\\
&=R i_!(\check E_U)(\deleted{\rho+}1)+\sum_{x\in Y}\deleted{\unsure{[}}j_{x*}[\check E_\eta^{\leftrightharpoons(x)}\deleted{\unsure{]|)}}+{}\\
&\hspace{21em}j_x^*\mathbb R i_{\unsure{*|x}}(\check E_U)(1)]\\
\text{\deleted{Il y a donc}}\quad &=\check E_\eta^{\leftrightharpoons}(1)+\sum_{x\in Y}j_{x*}(\mu_{\deleted{\unsure{x}}x}(E))
\end{aligned}
\]

\[
\begin{aligned}
\mu_x(E)&=\check E_\eta^{\leftrightharpoons(x)}\deleted{\unsure{\times|x|\otimes}}(\deleted{\unsure{i}}\unsure{Q_\ell|1_\ell|\mathbb{Q}_\ell}-Q_\ell(1))+j_x^*\mathbb R i_{\unsure{*\eta|\eta}}(\check E_\eta)(1)\\
&\deleted{=\check E_\eta^{\leftrightharpoons(x)}}
\end{aligned}
\]
Or \deleted{\unsure{par la du}} par dualité locale, on a
\[
j_x^*(\mathbb R i_\eta(\check E_\eta))=\check E_{\eta\deleted{\unsure{x|*}}}(x)-(E_\eta(x))^\vee(-1)
\]
% archive sheet number: 48
\scanpage{50}
\mspage{4 - 40 -}
n\textsuperscript{o} 392

\uline{indépendants sur} $k(\underline p)$ ; \uline{pour qu'ils soient linéairement indépendants sur} $k(\underline p)$ \uline{pour tout} $\underline p\in P$, \uline{il faut et il suffit que} $M/L$ \uline{soit sans torsion}.

Soient $x_{m+1},\ldots,x_n$ des éléments de $M$ qui, avec $x_1,\ldots,x_m$, forment une base de $V$, et soit $N$ le sous-$A$-module libre de $M$ engendré par les $x_i$, $(1\leq i\leq n)$. Il résulte du n\textsuperscript{o} 3, th. 3 que l'on a $N_{\underline p}=M_{\underline p}$ pour presque tout $\underline p\in P$ ; comme $x_1(\underline p),\ldots,x_n(\underline p)$ forment une base de $N(\underline p)$ sur $k(\underline p)$, cela établit la première assertion.

Si $M/L$ est sans torsion, il en est de même de $(M/L)_{\underline p}=M_{\underline p}/L_{\underline p}$ pour tout $\underline p\in P$, (n\textsuperscript{o} 1, \uline{Exemple} 6), et comme $A_{\underline p}$ est principal, $M_{\underline p}/L_{\underline p}$ est libre. Par suite, $M_{\underline p}$ est somme directe de $L_{\underline p}$ et d'un $A_{\underline p}$-module libre $E$ de rang $n-m$ ; donc $M(\underline p)$ est somme directe de $L(\underline p)$ et du $k(\underline p)$-espace vectoriel $E/\underline p E$, de dimension $n-m$ ; par suite $L(\underline p)$ est de dimension $m$, et comme il est engendré par les $x_i(\underline p)$, $(1\leq i\leq m)$, ces derniers sont linéairement indépendants.

Inversement, supposons que les $x_i(\underline p)$, $(1\leq i\leq m)$, soient linéairement indépendants sur $k(\underline p)$ pour tout $\underline p\in P$. Alors $L_{\underline p}$ est facteur direct de $M_{\underline p}$ pour tout $\underline p$, (chap. II, \S~3, n\textsuperscript{o} 2, cor. 1 et 2 de la prop. 5), et par suite $M_{\underline p}/L_{\underline p}=(M/L)_{\underline p}$ est sans torsion pour tout $\underline p\in P$. On en conclut que $P\handwritten{\cap}\operatorname{Ass}(M/L)=\varnothing$ en vertu du chap. IV, \S~1, n\textsuperscript{o} 2, cor. de la prop. 5. Mais comme $L$ est réflexif, il résulte du n\textsuperscript{o} 2, prop. 6, (i), que le seul idéal premier qui puisse appartenir à $\operatorname{Ass}(M/L)$ est l'idéal $(0)$ ; donc $M/L$ est sans torsion.

LEMME 8.- \uline{Supposons que le rang} $n$-\uline{de} $M$ \uline{soit} $\geq 2$ ; \uline{alors il existe un élément} $x\neq 0$ \uline{de} $M$ \uline{tel que} $M/Ax$ \uline{soit sans torsion}.

Soit $y\neq 0$ un élément de $M$. En vertu du lemme 7, l'ensemble $Y$ des $\underline p\in P$ tels que $y(\underline p)=0$ est fini. Si $Y=\varnothing$, il résulte du lemme 7 appliqué à la suite $(x_i)$ formée du seul élément\unsure{.} $y$, que $M/Ay$ est sans torsion. Supposons donc $Y\neq\varnothing$, et posons $S=\handwritten{\bigcap}_{\underline p\in Y}(A-\underline p)$ ; on sait, (n\textsuperscript{o} 4, lemme 2), que $S^{-1}A$ est un anneau
% archive sheet number: 49
\scanpage{51}
Posons
\[
\boxed{\alpha_x(E_\eta)=E_\eta^{\unsure{\leftrightharpoons|\leftrightarrows|\natural}(x)}-E_\eta(x)}
\]

Alors
\[
\boxed{\mu_x(E)=\alpha_x(E_\eta)^\vee-\alpha_x(\check E_\eta)(1)}
\]

$[$\unsure{donc|d'où} $\mu_x(\check E)=\alpha_x(\check E_\eta)^\vee-\alpha_x(E_\eta)(1)\deleted{\unsure{)}}=-\mu_x(E)^\vee(1)$ i.e.
\[
\begin{aligned}
&\deleted{\mu_x(\check E)\unsure{=-\mu_x|=\mu_x}}\\
&\left\{\begin{aligned}
\mu_x(\check E)&=-\mu_x(E)^\vee(1)\\
\mu_x(E)&=-\mu_x(\check E)^\vee(1)
\end{aligned}\right.
\end{aligned}\qquad ]
\]

Donc on \unsure{trouve}
\[
\boxed{\left\{\begin{aligned}
D(E_\eta^{\leftrightharpoons})&=\check E_\eta^{\leftrightharpoons}(1)+\sum_{x\in X^{(0)}}j_{x*}(\mu_x(E))\\
\mu_x(E)&=\alpha_x(E_\eta)^\vee-\alpha_x(\check E_\eta)(1)\qquad
\alpha_x(E_\eta)=E_\eta^{\leftrightharpoons(x)}-E_\eta(x)
\end{aligned}\right.}
\]

\uline{NB.}\quad $\alpha_x(E_\eta)$ \uline{n'est pas} additif en $E_\eta$, \uuline{mais} $\mu_x(E_\eta)$ \uuline{l'est}

Si $\check E_\eta=E_\eta(\rho)$, \unsure{cela donne|on trouve}
\[
\boxed{\begin{aligned}
D(E_\eta^{\leftrightharpoons})&=E_\eta^{\leftrightharpoons}(\rho+1)+\sum_{x\in X^{(0)}}j_{x*}(\mu_x(E)\\
\mu_x(E)&=\alpha_x(E_\eta)^\vee-\alpha_x(E_\eta)(\rho+1)
\end{aligned}}
\]
% archive sheet number: 50
\scanpage{52}
\margin{\handwritten{\unsure{1|\textbackslash}}}
\mspage{4 - 38 -}
n\textsuperscript{o} 392

où $I$ est pseudo-nul. Il résulte donc du n\textsuperscript{o} 7, prop. 16, (v) et cor. 1 de la prop. 16, que l'on a $c_B(M\otimes_A B)=c_B(Q\otimes_A B)=-c_B(\handwritten{\chi}_B(Q\otimes_A B))=-c_B(i(\handwritten{\chi}_A(Q)))$ en vertu de (ii) ; mais par définition de l'homomorphisme $i:C(A)\longrightarrow C(B)$ on a $c_B(i(\handwritten{\chi}_A(Q)))=i(c_A(\handwritten{\chi}_A(Q)))=-i(c_A(M))$, ce qui achève de prouver (23).

\uline{Remarques}.- 1) Si $M$ est un $A$-module réflexif, $M\otimes_A B$ n'est pas nécessairement réflexif, (exerc. \quad). Il en est toutefois ainsi lorsque $B$ est un $A$-module \uline{plat}, (n\textsuperscript{o} 2, prop. 8).

2) Soit $C$ un troisième anneau noethérien intégralement clos, tel que $B\subset C$ et que $C$ soit un $B$-module de type fini, (donc aussi un $A$-module de type fini). On a alors les formules de transitivité :
\[
N_{C/A}=N_{B/A}\circ N_{C/B}\ ,\tag*{\makebox[0pt][r]{\makebox[\displaywidth][l]{(25)}}}
\]
\[
i_{C/A}=i_{C/B}\circ i_{B/A}\tag*{\makebox[0pt][r]{\makebox[\displaywidth][l]{(26)}}}
\]
qui résultent immédiatement des formules de transitivité pour les indices de ramification et les degrés résiduels, (chap. VI, \S~8, n\textsuperscript{o} 1, lemme 1).

9. \uline{Un théorème de réduction}.

Les notations et hypothèses sont de nouveau celles des n\textsuperscript{os} 2 à 7.

LEMME 6.- \uline{Soient} $R$ \uline{un anneau commutatif}, $\underline p_i$, $(1\leq i\leq n)$, \uline{des idéaux premiers de} $R$, \uline{deux à deux distincts}.

(i) \uline{Pour} $1\leq i\leq n$, \uline{soit} $H_i$ \uline{une partie de} $R/\underline p_i$ \uline{vérifiant la condition suivante} : \uline{il n'existe aucun élément} $\alpha_i\in R/\underline p_i$ \uline{tel que} $\alpha_i+H_i$ \uline{contienne un idéal} $\neq 0$ \uline{de} $R/\underline p_i$. \uline{Alors il existe} $a\in R$ \uline{tel que}, \uline{pour} $1\leq i\leq n$, \uline{l'image canonique de} $a$ \uline{dans} $R/\underline p_i$ \uline{n'appartienne pas à} $H_i$.

(ii) \uline{Si} $\operatorname{Card}(H_i)<\operatorname{Card}(R/\underline p_i)$, \uline{les} $H_i$ \uline{vérifient la condition de} (i).

(i) Raisonnons par récurrence sur $n$, le cas $n=0$ étant trivial. Soit donc $n\geq 1$. Quitte à faire une permutation sur les indices $i$, on peut supposer que $\underline p_1$ est
% archive sheet number: 51
\scanpage{53}
L'équation \unsure{fonctionnelle} devient alors
\[
L_{\check E_\eta}(p^{-1}t)\prod_x\deleted{\unsure{\delta}}\lambda_x(E)=(-t)^{-\chi^{\unsure{\leftrightharpoons|\leftrightarrows|\natural}}(E_\eta)}\delta^{\leftrightharpoons}(E_\eta)L_{E_\eta}(t^{-1})
\]
\margin{i.e. \unsure{remplaçant} $t$ par $pt$}
% slanted

\[
\left\{\begin{aligned}
L_{\check E_\eta}\left(\frac1{pt}\right)&=\prod_x\lambda_x(E)\underbrace{(-pt)^{\chi^{\leftrightharpoons}(E_\eta)}\delta^{\leftrightharpoons}(E_\eta)^{-1}}L_{\check E_\eta}(t)\\
\deleted{\unsure{\alpha|a}}\chi^{\leftrightharpoons}(E_\eta)&=\chi(E_\eta^{\leftrightharpoons})\qquad \text{à déterminer ....}\\
\delta^{\leftrightharpoons}(E_\eta)&=\delta(E_\eta^{\leftrightharpoons})\qquad \text{à déterminer ....}\\
\lambda_x(E)&=L_{\mu_x(E)}(t)=L_{\alpha_x(E_\eta)^\vee}(t)/L_{\alpha_x(\check E_\eta)}(p^{-1}t)
=(-t)^{-\deleted{\unsure{\rho}}\chi_x(E_\eta)}\deleted{\unsure{\rho}}\delta_x(E_\eta)\frac{L_{\alpha_x(E_\eta)}(t^{-1})}{L_{\alpha_x(\check E_\eta)}(t)}
\end{aligned}\right.
\]

Supposons $E$ de poids $\rho$, donc
\[
\check E\simeq E(\rho)\ ,\qquad L_{\check E}(t)=L_E(p^{-\rho}t)\ ,
\]
il vient alors, en \unsure{remplaçant} $t$ par $p^\rho t$
\[
\boxed{L_{E_\eta}\left(\frac1{qt}\right)=\left(\prod_x\lambda_x(E)(t)\right)\delta^{\leftrightharpoons}(E_\eta)^{-1}(-\deleted{\unsure{p}}qt)^{\chi^{\leftrightharpoons}(E_\eta)}L_E(t)}
\]
\margin{\deleted{$q=p^{\unsure{1+\rho|\rho}}$}\quad \deleted{$\lambda_x$}}
% slanted
\[
\left[\begin{aligned}
\delta^{\leftrightharpoons}(E_\eta)^{-1}\deleted{\unsure{p}}q^{\chi^{\leftrightharpoons}(E_\eta)}
&\overset{\textcircled{?}}{=}\deleted{\unsure{q^{\chi}|\varepsilon(E)}}\delta^{\leftrightharpoons}(E_\eta)\\
\text{i.e.}\quad \delta^{\leftrightharpoons}(E_\eta)^2&=q^{\chi^{\leftrightharpoons}(E_\eta)}\ ??
\end{aligned}\right]\ \longleftarrow ??
\]

\[
\left\{\begin{aligned}
q&=p^{1+\rho}\\
\lambda_x(E_\eta)&=L_{\mu_x(E_\eta)}(t)\qquad
\mu_x(E_\eta)=\alpha_x(E_\eta)^\vee-\alpha_x(E_\eta)(\rho+1)\\
&\deleted{\unsure{\delta_x(E_\eta)t}}\vphantom{\frac11}\\
&\deleted{=\unsure{(-t)}}=(-t)^{-\deleted{\unsure{\rho|\beta}}\chi_x(E_\eta)}\delta_x(E_\eta)\frac{L_{\alpha_x(E_\eta)}(t^{-1})}{L_{\alpha_x(E_\eta)}(p^{-\rho}t)}\\
\deleted{\unsure{\rho|\beta}}\chi_x(E_\eta)&=\chi(\alpha_x(E_\eta))\\
\delta_x(E_\eta)&=\delta(\alpha_x(E_\eta))
\end{aligned}\right.
\]
% archive sheet number: 52
\scanpage{54}
\mspage{4 - 37 -}
n\textsuperscript{o} 392

lui aussi injectif. Si $I=\operatorname{Ker}(f\otimes 1)$, on a $I_{\underline p}=\operatorname{Ker}((f\otimes 1)_{\underline p})$, (chap. II, \S~2, n\textsuperscript{o} 4, th. 1) ; on a donc $I_{\underline p}=0$, d'où \uline{a fortiori} $I_{\underline P}=(I_{\underline p})_{\underline P}=0$ pour $\underline P\mid\underline p$, ce qui prouve le lemme, (n\textsuperscript{o} 4, prop. 9).

Revenons à la démonstration de (ii). Pour tout $A$-module de torsion de type fini $M$, posons $\varphi(M)=\handwritten{\chi}_B(M\otimes_A B)$ ; il résulte de (i) que si $M$ est pseudo-nul, on a $\varphi(M)=0$. D'autre part, considérons une suite exacte de $A$-modules de torsion de type fini :
\diagram{\begin{tikzcd}[ampersand replacement=\&,column sep=1.5em]
0\arrow[r] \& M_1\arrow[r] \& M_2\arrow[r] \& M_3\arrow[r] \& 0.
\end{tikzcd}}
Il résulte du lemme 4 que l'on a une suite exacte de $B$-modules :
\diagram{\begin{tikzcd}[ampersand replacement=\&,column sep=1.5em]
0\arrow[r] \& I\arrow[r] \& M_1\otimes_A B\arrow[r] \& M_2\otimes_A B\arrow[r] \& M_3\otimes_A B\arrow[r] \& 0
\end{tikzcd}}
où $I$ est pseudo-nul. En utilisant le n\textsuperscript{o} 5, cor. de la prop. 10, on a donc $\varphi(M_2)=\varphi(M_1)+\varphi(M_3)$. On conclut donc de la prop. 11 du n\textsuperscript{o} 5 qu'il existe un homomorphisme $\vartheta:D(A)\longrightarrow D(B)$ tel que $\varphi(M)=\vartheta(\handwritten{\chi}_A(M))$ pour tout $A$-module de torsion de type fini $M$. Pour prouver que $\vartheta=i$, il suffit de montrer que $\varphi(A/\underline p)=i(\underline p)$ pour tout $\underline p\in P(A)$ ; or on a $(A/\underline p)\otimes_A B=B/\underline p B$, et pour tout $\underline P\in P(B)$, $(B/\underline p B)_{\underline P}=B_{\underline P}/\underline p B_{\underline P}$ ; ce dernier module est 0 si $\underline P$ n'est pas au-dessus de $\underline p$ ; si au contraire $\underline P\mid\underline p$, $B_{\underline P}/\underline p B_{\underline P}$ est un $B_{\underline P}$-module de longueur $e(\underline P/\underline p)$ par définition de l'indice de ramification, (chap. VI, \S~8, n\textsuperscript{o} 1) ; on a donc $\handwritten{\chi}_B(B/\underline p B)=\sum_{\underline P\mid\underline p}e_{\underline P/\underline p}\unsure{\circ|\cdot}\underline P=i(\underline p)$, ce qui achève de prouver (ii).

La formule (24) est immédiate, car $(M\otimes_A B)\otimes_B L=M\otimes_A L=(M\otimes_A K)\otimes_K L$ et la dimension sur $L$ de $(M\otimes_A K)\otimes_K L$ est égale à la dimension sur $K$ de $M\otimes_A K$. Pour démontrer (23), considérons un sous-module libre $H$ de $M$ tel que $Q=M/H$ soit un $A$-module de torsion. Appliquant comme ci-dessus le lemme 4, on a une suite exacte\unsure{-}de $B$-modules :
\diagram{\begin{tikzcd}[ampersand replacement=\&,column sep=1.5em]
0\arrow[r] \& I\arrow[r] \& H\otimes_A B\arrow[r] \& M\otimes_A B\arrow[r] \& Q\otimes_A B\arrow[r] \& 0
\end{tikzcd}}
% archive sheet number: 53
\scanpage{55}
\[
L_{E_\eta}\left(\frac1{qt}\right)=A(t)L_{E_\eta}(t)
\]

\[
A(t)=\prod_x\tikz[remember picture,baseline=(localfactor.base)]{\node[draw,ellipse,inner sep=2pt] (localfactor) {$\displaystyle\frac{L_{\alpha_x(E_\eta)}(t^{-1})}{L_{\alpha_x(E_\eta)}(p^{-\rho}t)}$};}
(-t)^{\chi^{\unsure{\leftrightharpoons|\leftrightarrows|\natural}}(E_\eta)-\sum_x\chi_x(E_\eta)}
\frac{q^{\chi^{\leftrightharpoons}(E_\eta)}}{\delta^{\leftrightharpoons}(E_\eta)\prod_x\delta_x(E_\eta)^{-1}}
\]
\hspace{.42\linewidth}\tikz[remember picture,baseline=(factorlabel.base)]{\node[inner sep=0pt] (factorlabel) {$\inserted{=A_x(t^{-1})^{-1}}$};}
\tikz[remember picture,overlay]{\draw (localfactor.south east) to[out=-20,in=180] (factorlabel.west);}

\[
L_{\alpha_x(\deleted{\unsure{E}}E_\eta)}\left(\frac1{p^\rho t}\right)=A_x(t)L_{\alpha_x(E_\eta)}(t)
\]
\[
\deleted{L_{\alpha_x(\unsure{E_\eta})}\quad L_{\alpha_x(E_\eta)}(}
\]

N.B. Ayant coupé $\lambda_x(\deleted{\unsure{E_\eta}}E_\eta)(t)$ \deleted{\unsure{en}} en morceaux, on en perd le \uline{\unsure{caractère} additif}, Il vaut sans doute mieux ne pas faire ce découpage ---

\uline{Pb} a) Dans les cas géométriques \unsure{classiques}, construire les $\alpha_x(E_\eta)$ qui \unsure{déterminent} les $\mu_x$, donc les $\lambda_x(E_\eta)(t)$. Déterminer en particulier \tikz[remember picture,baseline=(chiloc.base)]{\node[inner sep=0pt] (chiloc) {$\chi_x(E_\eta)$};}\margin{\tikz[remember picture,baseline=(geolabel.base)]{\node[inner sep=0pt] (geolabel) {\shortstack{\unsure{invariant}\\\uuline{géom.}}};}}, \tikz[remember picture,baseline=(deltaloc.base)]{\node[inner sep=0pt] (deltaloc) {$\delta_x(E_\eta)$};}\margin{\tikz[remember picture,baseline=(arithlabel.base)]{\node[inner sep=0pt] (arithlabel) {\shortstack{\unsure{invariant}\\\uuline{arithm.}}};}}, les \unsure{valeurs absolues} des val propres de \unsure{frobenius|Frobenius}.

b) Donner une formule genre \unsure{Chafarevitch|Chafarévitch} - Ogg pour \tikz[remember picture,baseline=(chiglob.base)]{\node[inner sep=0pt] (chiglob) {$\chi^{\leftrightharpoons}(E_\eta)$};}

c) Donner une formule pour \tikz[remember picture,baseline=(deltaglob.base)]{\node[inner sep=0pt] (deltaglob) {$\delta^{\leftrightharpoons}(E_\eta)$};}
\tikz[remember picture,overlay]{
\draw[->] (chiloc.south) to[out=-85,in=155] (geolabel.west);
\draw[->] (chiglob.east) to[out=0,in=-105] (geolabel.south);
\draw[->] (deltaloc.south) to[out=-80,in=100] (arithlabel.north);
\draw[->] (deltaglob.east) to[out=0,in=-135] (arithlabel.south);
}
% archive sheet number: 54
\scanpage{56}
\mspage{4 - 36 -}
n\textsuperscript{o} 392

Revenant à la démonstration de la prop. 19, on a, en vertu du lemme 4, $\handwritten{\bigcap}_{\underline P\in P(B)}R_{\underline P}=\handwritten{\bigcap}_{\underline p\in P(A)}(R_{[A]})_{\underline p}$, et la conclusion résulte du n\textsuperscript{o} 2, th. 2.

PROPOSITION 20.- (i) \uline{Pour qu'un} $A$-\uline{module de type fini} $M$ \uline{soit pseudo-nul}, \uline{il faut et il suffit que} $M\otimes_A B$ \uline{soit un} $B$-\uline{module pseudo-nul}.

(ii) \uline{Si} $M$ \uline{est un} $A$-\uline{module de torsion de type fini}, $M\otimes_A B$ \uline{est un} $B$-\uline{module de torsion de type fini}, \uline{et l'on a} :
\[
\handwritten{\chi}_B(M\otimes_A B)=i(\handwritten{\chi}_A(M)).\tag*{\makebox[0pt][r]{\makebox[\displaywidth][l]{(22)}}}
\]

(iii) \uline{Si} $M$ \uline{est un} $A$-\uline{module de type fini}, $M\otimes_A B$ \uline{est un} $B$-\uline{module de type fini}, \uline{et l'on a} :
\[
c_B(M\otimes_A B)=i(c_A(M))\tag*{\makebox[0pt][r]{\makebox[\displaywidth][l]{(23)}}}
\]
\[
r_B(M\otimes_A B)=r_A(M).\tag*{\makebox[0pt][r]{\makebox[\displaywidth][l]{(24)}}}
\]

(i) Soient $\underline P$ un idéal premier de $B$, $\underline p=\underline P\handwritten{\cap} A$ ; on a $(M\otimes_A B)_{\underline P}=M\otimes_A B_{\underline P}$, (chap. II, \S~2, n\textsuperscript{o} 7, prop. 18), et d'autre part $M\otimes_A B_{\underline P}=(M\otimes_A A_{\underline p})\otimes_{A_{\underline p}}B_{\underline P}={}=M_{\underline p}\otimes_{A_{\underline p}}B_{\underline P}$ ; la relation $M_{\underline p}=0$ est donc équivalente à $(M\otimes_A B)_{\underline P}=0$, (chap. II, \S~4, n\textsuperscript{o} 4, lemme 4). Il suffit d'appliquer cette remarque à l'idéal $\underline P=(0)$ et aux idéaux $\underline P\in P(B)$ pour prouver (i), compte tenu du n\textsuperscript{o} 4, prop. 9.

Pour prouver (ii), nous utiliserons le lemme suivant :

LEMME 5.- \uline{Soient} $M_1$, $M_2$ \uline{deux} $A$-\uline{modules de type fini}, $f:M_1\longrightarrow M_2$ \uline{un homomorphisme injectif}. \uline{Alors le noyau de} $f\otimes 1_B:M_1\otimes_A B\longrightarrow M_2\otimes_A B$ \uline{est pseudo-nul}.

Soit $\underline p$ un idéal premier de $A$ de hauteur $\leq 1$. On a $(M_i\otimes_A B)_{\underline p}=(M_i)_{\underline p}\otimes_{A_{\underline p}}B_{\underline p}$ $(i=1,2)$, (chap. II, \S~2, n\textsuperscript{o} 7, prop. 18), et $(f\otimes 1_B)_{\underline p}=f_{\underline p}\otimes 1_{B_{\underline p}}$ ; l'hypothèse que $f$ est injectif entraîne qu'il en est de même de $f_{\underline p}$, (chap. II, \S~1, n\textsuperscript{o} 4, th. 1) ; d'autre part, vu le choix de $\underline p$, $A_{\underline p}$ est un anneau principal et $B_{\underline p}$ un $A_{\underline p}$-module sans torsion de type fini, donc \uline{libre} ; on en conclut que $f_{\underline p}\otimes 1_{B_{\underline p}}$ est
% archive sheet number: 55
\scanpage{57}
Sur la formule virtuelle \unsure{donnant} $D(E_\eta^{\unsure{\leftrightharpoons|\leftrightarrows|\natural}})$, on lit
\[
\delta^{\leftrightharpoons}(E_\eta)^{-1}=\delta^{\leftrightharpoons}(E_\eta)p^{-\chi^{\leftrightharpoons}(E)(\rho+1)}\prod_x\delta_x(\mu_x(E))
\]

Or
\[
\begin{aligned}
\delta(\mu_x(E))&=\delta(\alpha_x(E_\eta))^{-1}\big[\delta(\alpha_x(\unsure{E_\eta|\eta}))p^{-\chi_x(E_{\unsure{\eta|\rho}})(\rho+1)}\big]^{-1}\\
&=\delta_x(E_\eta)^{-2}q^{+\chi_x(E_\eta)}
\end{aligned}
\]

d'où
\[
\deleted{\unsure{(|\lbrace}}\prod_x\frac{\delta_x(E_\eta)^2}{q^{\chi_x(E_\eta)}}\deleted{\unsure{)|\rbrace}}=
\delta^{\leftrightharpoons}(E_\eta)^2q^{\sum\chi_x(E_\eta)}
\]

\[
\left\{\begin{aligned}
\frac{\delta^{\leftrightharpoons}(E_\eta)^2}{q^{\chi^{\leftrightharpoons}(E)}}\deleted{=\,q^{\chi^{\leftrightharpoons}(E)}}&=\prod_x\frac{\delta_x(E_\eta)^2}{q^{\chi_x(E_\eta)}}\qquad \text{\unsure{ou encore|on aura}}\\
\Vert\qquad&\\
\varepsilon^{\leftrightharpoons}(E_\eta)^2&=\prod_x\frac{\delta_x(E_\eta)^2}{q^{\chi_x(E_\eta)}}
\end{aligned}\right.
\]

\[
\begin{aligned}
\varepsilon^{\leftrightharpoons}(E_\eta)&=\pm\prod_x\frac{\delta_x(E_\eta)}{q^{\chi_x(E_\eta)/2}}
=\left(\prod_x\frac{\delta_x(E_\eta)}{p^{\rho\chi_x(E_\eta)/2}}\right)\tikz[remember picture,baseline=(upperpower.base)]{\node[inner sep=0pt] (upperpower) {$\displaystyle\frac{1}{p^{\chi/2}}$};}\\
\Vert\qquad&\\
\frac{\delta^{\leftrightharpoons}(E_\eta)}{q^{\deleted{\unsure{\delta}}\chi^{\leftrightharpoons}(E_\eta)/2}}&\qquad\qquad\qquad\tikz[remember picture,baseline=(sumchi.base)]{\node[inner sep=0pt] (sumchi) {$\inserted{\chi=\sum_x\chi_x(E_\eta)}$};}
\end{aligned}
\]
\[
\varepsilon^{\leftrightharpoons}(E_\eta)=\pm\left(\prod_x\varepsilon_x(E_\eta)\right)\tikz[remember picture,baseline=(lowerpower.base)]{\node[inner sep=0pt] (lowerpower) {$p^{-\chi/2}$};}
\]
\tikz[remember picture,overlay]{
\draw[->] (sumchi.north west) to[out=10,in=-110] (upperpower.south west);
\draw[->] (sumchi.south east) to[out=-160,in=75] (lowerpower.north east);
}
% archive sheet number: 56
\scanpage{58}
\mspage{4 - 35 -}
n\textsuperscript{o} 392

libre de $R$ tel que $R/M$ soit un $B$-module de torsion ; comme $M_{[A]}$ est somme directe de $r_B(R)$ $A$-modules isomorphes à $B$, on a, (prop. 16, (i)), $c_A(M_{[A]})={}=r_B(R)\cdot c_A(B)$. D'autre part, $c_A((R/M)_{[A]})=-c_A(N(\handwritten{\chi}_B(R/M)))$ en vertu de (19) ; mais par définition de l'homomorphisme $N:C(B)\longrightarrow C(A)$, on a $c_A(N(d))=N(c_B(d))$ pour tout $d\in D(B)$, et comme $c_B(\handwritten{\chi}(R/M))=-c_B(R)$ par définition, on a finalement $c_A((R/M)_{[A]})=N(c_B(R))$ ; il suffit alors d'appliquer la prop. 16, (i) pour obtenir (20).

PROPOSITION 19.- \uline{Soit} $R$ \uline{un} $B$-\uline{module de type fini}. \uline{Pour que} $R$ \uline{soit réflexif}, \uline{il faut et il suffit que} $R_{[A]}$ \uline{soit un} $A$-\uline{module réflexif}.

Il est clair que pour que $R$ soit un $B$-module sans torsion, il faut et il suffit que $R_{[A]}$ soit un $A$-module sans torsion. On peut donc supposer que $R$ est un réseau de $W=R\otimes_B L$ par rapport à $B$. Nous ut\deleted{\unsure{l|i}}iliserons le lemme suivant :

LEMME 4.- \uline{Soit} $W$ \uline{un espace vectoriel de dimension finie sur} $L$ \uline{et soit} $R$ \uline{un réseau de} $W$ \uline{par rapport à} $B$. \uline{Pour tout} $\underline p\in P(A)$, \uline{on a alors} $(R_{[A]})_{\underline p}=\handwritten{\bigcap}_{\underline P\mid\underline p}R_{\underline P}$.

En effet, si $S=A-\underline p$, les idéaux premiers de l'anneau $S^{-1}B$ sont engendrés par les idéaux premiers de $B$ ne rencontrant pas $S$, autrement dit les idéaux $\underline P_i$, $(1\leq i\leq m)$, au-dessus de $\underline p$ et l'idéal $(0)$ ; cela montre que $S^{-1}B$ est un anneau semi-local dont les idéaux maximaux sont les $\underline m_i=\underline P_i(S^{-1}B)$ pour $1\leq i\leq m$ ; en outre l'anneau local $(S^{-1}B)_{\underline m_i}$ est isomorphe à $B_{\underline P_i}$, (chap. II, \S~2, n\textsuperscript{o} 5, prop. 11)\unsure{,} donc est un anneau de valuation discrète. L'anneau $S^{-1}B$ est par suite un anneau de Dedekind, ( \S~3, n\textsuperscript{o} 2, th. 1, e)), et comme il est semi-local, il est principal, ( \S~3, n\textsuperscript{o} 2, prop. 1). Cela étant, $(R_{[A]})_{\underline p}$ est égal à $S^{-1}R$, considéré comme $A_{\underline p}$-module ; d'après ce qui précède, $S^{-1}R$ est un réseau \uline{libre} de $W$ par rapport à $S^{-1}B$ et on peut par suite lui appliquer le th. 2 du n\textsuperscript{o} 2, qui donne $S^{-1}R=\handwritten{\bigcap}_i(S^{-1}R)_{\underline m_i}$ ; mais $(S^{-1}R)_{\underline m_i}=R_{\underline P_i}$, ce qui prouve le lemme.
% archive sheet number: 57
\scanpage{59}
On voit donc que c) est \unsure{ramené} (\uline{modulo un signe}) au cas de b), donc à déterminer $\chi^{\unsure{\leftrightharpoons|\leftrightarrows|\natural}}(E_\eta)$. Or cette \unsure{détermination} est faite, en utilisant les représentations d'Artin -{}-{}-{}-{}-
% archive sheet number: 58
